

%%%%%%%%%%%%%%%
\begin{frame}[fragile]{Transactions sur la base Messagerie}

\begin{minted}[baselinestretch=.9,
               fontsize=\small,
               framesep=2mm]{python}
# Tous les messages non envoyés
messages = connexion.cursor()
messages.execute("select * from Message where dateEnvoi is null")
for message in messages.fetchall():
    # Marquage du message
    connexion.begin()
    maj = connexion.cursor()
    maj.execute ("Update Message set dateEnvoi='2018-12-31' "
                + "where idMessage=%s", message['idMessage'])
                
    # Ici on envoie les messages à tous les destinataires
               
    # Ici on valide
    connexion.commit()
\end{minted}

\vfill

\alert{Est-elle à l'abri de tout problème de concurrence\,?}
\end{frame}


%%%%%%%%%%%%%%%
\begin{frame}{Analyse de la transaction}

On crée une transaction par message. Chaque transaction:

\vfill

\begin{redbullet}
  \item Lit un message sans date d'envoi dans le curseur
   \item Modifie le message dans la base en affectant la date d'envoi
  \item Envoie le message
   \item Valide (\sqlkw{commit})
\end{redbullet}

\vfill

\begin{blocgauche}
\alert{Si chaque transaction s'exécutait seule et sans erreur,
chaque message serait envoyé exactement une fois.}
\end{blocgauche}
\vfill

Regardons de plus près...

\end{frame}

%%%%%%%%%%%%%%%
\begin{frame}{Que se passe-t-il en cas de panne\,?}

Imaginons (pire scénario) une panne \alert{juste avant} le \sqlkw{commit}
\vfill

\figSlideWithSize{transaction-messages-1.pdf}{15cm}
\vfill

\only<2-2>{

\begin{blocgauche}
\begin{redbullet}
   \item Le système effectue un \sqlkw{rollback}\,: le message se retrouve sans date d'envoi (à \sqlkw{null})
   \item Le message a quand même été envoyé (le SGBD n'a aucun contrôle là-dessus).
\end{redbullet}

\end{blocgauche}
}

\end{frame}

%%%%%%%%%%%%%%%
\begin{frame}{Isolation\,: et si une autre transaction modifie le message}

Deux exécutions concurrentes du programme d'envoi
\vfill

\figSlideWithSize{transaction-messages-2.pdf}{15cm}
\vfill

\only<2-2>{

\begin{blocgauche}
\begin{redbullet}
   \item Le résultat du curseur est constitué au moment du \texttt{execute}
   \item On continuera à lire des messages même s'ils ont été modifiés ou détruits entretemps\,!
\end{redbullet}

\end{blocgauche}
}
\end{frame}


%%%%%%%%%%%%%%%
\begin{frame}{Isolation et exécutions simultanées}

Deux scripts d'envoi de message sont envoyés \alert{presque} en même temps

\vfill

\figSlideWithSize{transaction-messages-3.pdf}{14cm}
\vfill

\only<2-2>{

\begin{blocgauche}
\begin{redbullet}
   \item \alert{Isolation totale} (mode \sqlkw{serializable})\,: on aboutira à des \textit{deadlocks}
   (les lectures bloquent les écritures)
   \item Sinon (autres modes), pas de blocage, et défaut d'isolation\,:  les messages seront envoyés deux fois...
\end{redbullet}
}

\end{blocgauche}
\end{frame}

\begin{frame}{À retenir}

Toujours utile de se poser la question ``Comment ma transaction se comporte-t-elle
en cas de concurrence''
\vfill

\begin{redbullet}
  \item En cas de doute il est préférable de se mettre en mode \texttt{serializable}
  \item Les requêtes fournissant des données à une transaction doivent être \alert{dans} 
        la transaction.
\end{redbullet}

\vfill

\begin{blocgauche}
On pourrait envoyer de l'argent au lieu de simples messages, et la
question deviendrait très sensible...

\end{blocgauche}
\end{frame}

